\frfilename{mt27.tex} 
\versiondate{26.8.13} 
\copyrightdate{1996} 
 
\loadeusm 
\def\hatcalBnu{\hat{\Cal B}_{\nu}} 
 
\def\chaptername{Teori probabilitas} 
\def\sectionname{Pendahuluan} 
 
\newchapter{27} 
 
Lebesgue menciptakan teori integrasinya sebagai tanggapan terhadap sejumlah
masalah dalam analisis real, dan sepanjang hidupnya tampaknya ia memandang
teori tersebut sebagai alat untuk digunakan dalam geometri dan kalkulus
({\smc Lebesgue 72}, jilid\ 1 dan 2). Yang luar biasa, setelah disesuaikan
sebagaimana mestinya, teori itu ternyata menyediakan landasan yang kokoh bagi
teori probabilitas. Pengembangan pendekatan ini umumnya dikaitkan dengan nama
Kolmogorov. Pendekatan tersebut telah begitu mendominasi teori probabilitas
abstrak modern sehingga banyak penulis mengabaikan semua metode lain. Saya
tidak bermaksud mengikat diri pada pandangan apa pun tentang apakah ukuran
$\sigma$-aditif merupakan satu-satunya cara untuk memberi landasan yang ketat
bagi teori probabilitas, atau apakah ukuran-ukuran itu memadai untuk menangani
semua gagasan probabilistik; ada beberapa persoalan filosofis yang serius di
sini, karena teori probabilitas, setidaknya dalam aspek terapannya, berupaya
membantu kita memahami dunia material di luar matematika. Namun dari posisi
saya sebagai seorang ahli teori ukuran, tidak dapat disangkal bahwa teori
probabilitas termasuk di antara penerapan utama konsep dan teorema teori
ukuran, serta merupakan salah satu sumber gagasan baru yang paling hidup; dan
bahwa setiap ahli teori ukuran harus tanggap terhadap intuisi yang dapat
diberikan oleh metode probabilistik. 
 
Saya telah menulis paragraf sebelumnya dengan ungkapan yang menyiratkan bahwa
`teori probabilitas' dengan suatu cara dapat dibedakan dari bagian lain teori
ukuran; ini adalah perkara lain yang mengenainya saya lebih suka tidak
mengemukakan pendapat apa pun sebagai sesuatu yang definitif. Namun tidak
diragukan lagi terdapat suatu perbedaan, yang lebih dalam daripada fakta
elementer bahwa probabilitas (hampir) secara eksklusif menangani ruang
berukuran $1$.
M.Lo\`eve mengemukakan argumen yang meyakinkan ({\smc Lo\`eve 77}, $\S$10.2)
bahwa hakikat teori probabilitas terletak pada sifat artifisial ruang
probabilitas itu sendiri. Dalam teori ukuran, ketika kita ingin
mengintegralkan suatu fungsi, biasanya kita merasa bahwa kita mempunyai fungsi
yang sungguh-sungguh dengan domain dan nilai. Dalam teori probabilitas, ketika
kita mengambil ekspektasi suatu variabel acak, variabel tersebut merupakan suatu
`yang dapat diamati' atau `hasil suatu percobaan'; pada umumnya kita tidak
pasti, tidak mengetahui, atau tidak memedulikan faktor-faktor yang mendasari
variabel tersebut. Izinkan saya memberikan contoh dari teorema-teorema di bawah
ini. Dalam bukti Teorema Limit Pusat (274F), saya mendapati bahwa saya
memerlukan daftar bantu $Z_0,\ldots,Z_n$ berupa variabel-variabel acak yang saling
bebas dan bebas dari barisan asal $X_0,\ldots,X_n$. Saya menciptakan
barisan semacam itu dengan mengambil ruang produk $\Omega\times\Omega'$, dan
menulis $X'_i(\omega,\omega')=X_i(\omega)$, sedangkan $Z_i$ merupakan fungsi
dari $\omega'$. Sekarang perbedaan antara $X_i$ dan $X'_i$ adalah perbedaan
yang biasanya akan dianggap serius oleh seorang analis yang terlatih. Kita
tidak menganggap fungsi $x\mapsto x^2:[0,1]\to[0,1]$ sebagai hal yang sama
dengan fungsi $(x_1,x_2)\mapsto x_1^2:[0,1]^2\to[0,1]$. Namun seorang ahli
probabilitas kemungkinan merasa bahwa mulai menulis $X'_i$ sebagai ganti
$X_i$ benar-benar terlalu rewel. Sejak awal ia tidak memercayai ruang
$\Omega$, dan jika ruang itu ternyata tidak memadai bagi intuisinya, ia
memperluasnya tanpa ragu. Lo\`eve menyebut ruang probabilitas sebagai `fiksi',
`ciptaan imajinasi' dalam ungkapan Larousse; ruang-ruang itu diperlukan dalam
model yang diajarkan Kolmogorov kepada kita untuk digunakan, tetapi kita
mempunyai kebebasan yang sangat besar dalam memilihnya, dan pada hakikatnya
ruang-ruang tersebut sama sekali tidak setegas suatu himpunan dengan
titik-titik. 
 
Karena itu, suatu ruang probabilitas merupakan entitas yang agak lebih samar
dalam teori probabilitas daripada dalam teori ukuran. Objek-objek penting
dalam teori probabilitas adalah variabel acak dan distribusi, khususnya distribusi
gabungan. Dalam jilid ini saya akan menangani secara eksklusif variabel acak yang
dapat dipandang bernilai dalam suatu pangkat Kartesius dari $\Bbb R$; tetapi ini bukanlah
pokok utamanya. Yang penting ialah bahwa dengan suatu cara {\it kodomain},
himpunan nilai potensial, dari suatu variabel acak jauh lebih terdefinisi dengan
baik daripada {\it domain}-nya. Akibatnya perhatian kita tidak dipusatkan pada
ciri apa pun dari ruang artifisial yang nyaman digunakan sebagai ruang
probabilitas yang mendasari -- saya menulis `mendasari', meskipun itu merupakan
aspek model yang paling dangkal dan paling mudah diubah -- melainkan pada
distribusi di kodomain yang diinduksi oleh variabel acak. Jadi Teorema Limit
Pusat, yang hanya berbicara tentang distribusi, sebenarnya lebih penting dalam
probabilitas terapan daripada Hukum Kuat Bilangan Besar, yang menyatakan nilai
yang hampir pasti dicapai oleh suatu rata-rata jangka panjang. 
 
W.Feller ({\smc Feller 66}) bahkan melangkah lebih jauh daripada Lo\`eve, dan
sejauh mungkin bekerja sepenuhnya dengan distribusi, dengan menyusun perangkat
yang memungkinkannya menempuh bagian-bagian panjang tanpa menyebut ruang
probabilitas sama sekali. Saya tidak berusaha menirunya. Namun pendekatan
tersebut memberi pelajaran dan setia pada hakikat pokok bahasannya. 
 
Teori probabilitas mencakup lebih banyak matematika daripada yang dapat dengan
mudah dikuasai dalam satu masa hidup, dan untuk bab pengantar ini saya telah
memilih dua teorema limit yang sudah saya sebutkan, Hukum Kuat Bilangan Besar
dan Teorema Limit Pusat, bersama dengan sejumlah materi tentang martingal
(\S\S275-276). Keduanya menggambarkan bukan hanya watak khusus teori
probabilitas -- sehingga Anda dapat membentuk penilaian sendiri atas komentar
di atas -- melainkan juga beberapa sumbangan utamanya bagi teori ukuran
`murni', yaitu konsep `kebebasan' dan `ekspektasi bersyarat'. 
 
 
\discrpage 
 
